% Generated draft: requires mathematical and source review.
\documentclass[11pt]{article}
\usepackage[margin=25mm]{geometry}
\usepackage{fontspec}
\setmainfont{Latin Modern Roman}
\usepackage{amsmath,amssymb,mathtools,unicode-math}
\setmathfont{Latin Modern Math}
\usepackage{enumitem,needspace,xurl,hyperref}
\setlength{\parindent}{0pt}
\setlength{\parskip}{5pt}
\setlength{\emergencystretch}{3em}
\newcommand{\N}{\mathbb{N}}
\newcommand{\Z}{\mathbb{Z}}
\newcommand{\Q}{\mathbb{Q}}
\newcommand{\R}{\mathbb{R}}
\newcommand{\C}{\mathbb{C}}
\newcommand{\F}{\mathbb{F}}
\newcommand{\A}{\mathbb{A}}
\newcommand{\PP}{\mathbb{P}}
\newcommand{\E}{\mathbb{E}}
\newcommand{\eps}{\varepsilon}
\newcommand{\dd}{\,\mathrm{d}}
\begin{document}
\section*{UnsolvedMath ID: 30000996}
\textbf{Source problem code:} catalogue entry
\subsection{Definitions and mathematical statement}
The standard Heisenberg group \(\mathbb H^m\) is \(\R^m\times\R^m\times\R\) with product
\[
(x,y,t)(x',y',t')=(x+x',y+y',t+t'+\tfrac12(x\cdot y'-y\cdot x')),
\]
usually equipped with its left-invariant horizontal distribution and sub-Riemannian metric. The source does not specify \(m\), the metric normalization, or whether its area is horizontal perimeter or another related surface functional.

For the area functional \(\mathcal A\) and admissible surface class intended in the report, a critical point is a surface \(\Sigma\) for which the appropriate first variation vanishes,
\[
\left.\frac{\dd}{\dd s}\mathcal A(\Sigma_s)\right|_{s=0}=0,
\]
for all admissible compactly supported variations \((\Sigma_s)\) through \(\Sigma\), whenever this derivative is defined. The question is: what regularity follows for such area-critical surfaces, and under what natural hypotheses are they smooth? The parenthetical term ``minimal surfaces'' is retained in this stationarity sense; it does not silently impose the stronger property of minimizing area among all competitors.

The report extraction does not give the area formula, weak or classical notion of criticality, admissible variations, initial regularity, graph or embedding hypotheses, boundary conditions, treatment of characteristic points, or Euclidean versus intrinsic meaning of smoothness. For horizontal surface problems, characteristic points are points where the horizontal distribution is tangent to the surface; behavior there may need a separate formulation. No universal smoothness assertion, regularity exponent, or singular-set estimate is specified by the record, and none is invented.
\subsection{Short English statement}
What regularity do area-stationary surfaces in the Heisenberg group have, and when must they be smooth? The area functional, surface class, and intended hypotheses are not specified.
\subsection{Sources}
\begin{itemize}
\item Local UnsolvedMath registry, permanent ID \(30000996\), source code catalogue entry, ``Regularity of Minimal Surfaces in the Heisenberg Group.'' No individual proposer is supplied.
\item \emph{Calculus of Variations} (2009), Oberwolfach Reports, DOI 10.4171/owr/2008/31, pp. 1707--1770: \url{https://doi.org/10.4171/owr/2008/31}.
\item Supplied reference, ``partial regularity paper'': \url{https://doi.org/10.1016/j.na.2015.03.013}.
\item Supplied reference, ``recent Heisenberg minimal-surface work'': \url{https://doi.org/10.1515/agms-2025-0021}.
\item Supplied literature triage, checked 2026-08-21: status ``open.'' It reports that the source calls general regularity of area-critical points open and that subsequent work gives partial regularity without a general resolution. This dated assessment derives from supplied OWR curation; time-sensitive status and cited claims require independent verification.
\end{itemize}
\end{document}
